\backhead{FURTHER EXERCISES}{khred}
\begin{exlist}
\item A maize crop has $K_{c,ini}=0.3$, $K_{c,mid}=1.2$, $K_{c,end}=0.6$, with stages of 25, 40, 45 and 30 days. Find $K_c$ on days 10, 45, 90 and 125.\exans{0.30, 0.75, 1.20, 0.90}
\item With $ET_0=5.5$ mm\,d$^{-1}$ on day 90 of the crop above, $\TAW=140$ mm, $p=0.55$ and a depletion of 100 mm, find $K_s$ and the actual water use.\exans{0.635; 4.19 mm\,d$^{-1}$}
\item Over a season a crop uses 420 mm out of a potential 560 mm. With $K_y=1.25$ and $Y_m=11$ Mg\,ha$^{-1}$, find the expected yield and the water productivity of the actual water use.\exans{7.56 Mg\,ha$^{-1}$; 1.80 kg\,m$^{-3}$}
\item The evaporation from a bare soil in stage 2 follows $E=\beta\sqrt t$ with $\beta=4$ mm\,d$^{-1/2}$. How much evaporates in the first 4 days and in the next 5 days of stage 2?\exans{8 mm; 4 mm}
\item A crop with $K_c=1.05$ on a day with $ET_0=6.2$ mm\,d$^{-1}$ is irrigated by a drip system that wets 40 per cent of the surface. If soil evaporation is proportional to the wetted fraction and is 15 per cent of $ET_c$ under full wetting, how much is saved per day?\exans{0.59 mm}
\end{exlist}
\backhead{ESSAY QUESTIONS}{khred}
{\small\itshape Answer each question in about one page, with equations, a sketch and at least one number from this chapter.}\par\vspace{3pt}
\begin{exlist}
\item Explain the physical basis of the Penman--Monteith equation and why it is preferred to empirical equations for irrigation scheduling.
\item Discuss the single and the dual crop coefficient methods. When does the separation of soil evaporation from transpiration matter?
\item Explain how water stress reduces transpiration and yield. Why is the response of yield to water not linear for all crops?
\item Discuss the crop library of SWAFL: what information each crop carries and how a change of crop changes the other layers.
\item Compare the crop coefficients of a crop grown under drip and under sprinkler irrigation, and explain the difference.
\item Why is evaporation from the soil surface a target of water saving, and what are the limits of such savings?
\end{exlist}
\dykband{EXERCISES IN SWAFL}
{\small Open SWAFL (desktop or online) with the benchmark project. The chapter of the \href{https://github.com/halilunaziru73-creator/Coupled-Water-Fluxes-Pore-Scale-Dynamics-Precision-Irrigation-Mulching-Feedback-Loops/raw/main/Software/v1.1.0/SWAFL_Desktop/docs/manual_v1.1.0/SWAFL_Manual_v1.1.0.pdf}{SWAFL manual} that explains each task is given in brackets; record your values in a table and compare them with what is expected (\emph{Check}).}\par\vspace{3pt}
\begin{exlist}
\item {\color{khblue}[Manual, Chapter 21]} Open \emph{Crop water requirement} and choose maize, then wheat. Compare the $K_c$ curves and the seasonal water need.\par\noindent{\small\itshape Check: maize has a higher mid-season $K_c$ and a later peak.}
\item {\color{khblue}[Manual, Chapter 21]} In the same layer change the planting date by one month and record the change of the seasonal requirement.\par\noindent{\small\itshape Check: a later planting shifts the peak into the hotter period.}
\item {\color{khblue}[Manual, Chapter 19]} Import a weather file or use \emph{The weather of a place} for your own town and compare $ET_0$ with the benchmark record.\par\noindent{\small\itshape Check: the reference demand reflects your climate.}
\item {\color{khblue}[Manual, Chapter 22]} Open \emph{Reliability over 20 seasons} and read the years in which the crop was stressed. Relate them to the weather of those years.\par\noindent{\small\itshape Check: stress years coincide with dry, hot seasons.}
\end{exlist}
